\chapter{51}

Nach einigen Hundert Metern trennten sich ihre Wege. Mazi verabschiedete
sich in Richtung seiner Sennhütte. Flexa machte jede seiner Bewegungen
synchron mit. Vielleicht war es diese selbstverständliche Ordnung, die
Louis faszinierte. Diese gegenseitige Abstimmung, die ihn an seine
verlorene Sicherheit erinnerte.

Mit Manuela und Joh erreichte Louis endlich das kleine Holzhaus, das am
Berg zu kleben schien. In den leeren Raum gestellt, wirkte es, als warte
es geduldig auf seine Bestimmung.

Louis befreite sich mit einem Ächzen von all dem umgehängten Gepäck. Er
freute sich auf die freien Stunden an diesem menschenleeren Ort. Ganz
mit sich selbst. Und ohne beobachtende Blicke, die jedes Mal Irritation
auslösten.

Manuela hatte recht, die Hütte war einfach gehalten und wirkte kleiner,
als er sie sich in seiner Vorstellung ausgedacht hatte. Die Fotos im
Instruktionsordner mussten im Weitwinkel aufgenommen worden sein. Und
doch war ein spezielles Flair auszumachen. Eine Gemütlichkeit, die sich
von der Umgebung direkt ins Haus übertrug. Da war aber auch diese schwer
auszumachende Melancholie. Er hatte sie zwangsläufig mit hochgetragen.
Wahrgenommen hatte er sie beim Auspacken des Rucksacks. Als sie ihn
urplötzlich streifte. Wie ein luftiges Gespenst, das sich so rasch wie
gekommen, wieder davon machte. Er konnte es kaum erwarten, sich
musikalisch aufzutanken und seine Bürde vorübergehend abzulegen.

Joh und Manuela wollten am späten Nachmittag absteigen. Er würde die
Freiluftdusche einweihen und versuchen, seine Getriebenheit in den
Pausenmodus zu verbannen. Der Plan beflügelte ihn.

Bevor er die steile Leiter in den Raum hinunter stieg, kreisten seine
Gedanken um mögliche Musikstücke. Erhellen sollten sie ihn. Schon die
Vorstellung, sich endlich wieder einmal Melodien, mitreissenden Rhythmen
und einer wohltuenden Stimme hinzugeben, ganz für sich, löste ein
behagliches Gefühl aus. \textbf{Álvaro Soler}, den brauchte er jetzt.
Die letzten Monate hatte sich Louis seine Musik nicht mehr angehört.
Louis’ Leichtigkeit war bereits vor Giorgias Trennung zu Mangelware
geworden.

Sein Blick verlor sich in den Ritzen der Holzwand. Was war er für ein
armer Kerl. Wie konnte sein Leben nach dieser Geschichte überhaupt
weitergehen?

Manuela fand er unten in der Küchenecke verschiedene Nahrungsmittel
einräumen. Durch das kleine Südfenster sah er Joh. Vornüber gebeugt und
mit blossem Oberkörper löste dieser gerade die Schnürsenkel seiner
Bergschuhe. Als könne er eine Abkühlung kaum erwarten, stieg Joh
blitzschnell aus seinen restlichen Kleidern. Schliesslich stand er nackt
im Brunnentrog und spritze sich Wasser die Beine hoch. Der Anblick
erhellte Louis. Manuela stand, in der Hand eine Packung Spaghetti,
hinter ihm und kicherte.

«Joh live, der lässt sich nicht aufhalten,» sie klang vergnügt, «egal
wie wenig Wasser da drin ist.»

Louis drehte sich zu ihr um und lachte.

«Bist du bereit für ein paar Instruktionen, Ciro?»

Sie riss die Spaghetti-Packung auf, liess den Inhalt in das kochende
Wasser auf dem Herd gleiten und stützte ihre Arme in den Hüften ab. Sie
schien sich zu sammeln.

«Sicher.»

Louis trat näher an Manuela heran. Ihre ersten Informationen drehten
sich um Einzelheiten rund ums Kochen. Sie öffnete die kleinen Türchen
eines hohen Schrankes. Darin befanden sich Geschirr, Besteck und
sämtliche weiteren Küchengeräte. Das professionelle Messersortiment fiel
Louis sofort auf. Der Anblick trug ihn unweigerlich in die Küche des
Ristorante zurück. Er liess sich nichts anmerken.

Dann kam Joh, sichtlich erfrischt und in einer kurzen Hose, in die Hütte
getreten. Er trug einen Beutel mit verschiedenen Gemüsen mit sich,
leerte diesen auf dem Tisch aus und entnahm einer kleinen Schublade ein
Messer.

«Den Käse hast du eingepackt, Manuela, ja?» fragte er zwischen zwei
Schnitten und blickte zu ihr hoch.

«Ja, alles da,» sie zeigte auf die Kühlbox, «ich geh mit Ciro mal einen
Alptag durch, okay? Kommst du zurecht?»

Manuela strich Joh über den blossen Rücken. Er hob den Daumen, rückte
seine Brille zurecht und schnippelte weiter eine
\href{https://www.gemuese.ch/gemuesearten/zucchetti}{Zucchetti in kleine
Stücke.}

Manuelas Instruktionsidee war nicht nur informativ, sondern auch
erheiternd. Sie spielte mit Louis einen durchschnittlichen Alptag durch,
wie sie es nannte. Begonnen beim Aufstehen bis zum Zubettgehen am Abend.
Sie kletterten die Leiter hoch, hinunter, zirkelten unten um den
kochenden Joh herum, traten vor die Hütte, umrundeten sie und blieben
final bei der Freiluftdusche stehen. Louis’ Fragen stellten sich von
alleine ein.

«Hier also unser Highlight, schau, Ciro,» Manuela wies mit dem Finger
auf den langen, schlanken Tank, «der ist echt schnell leer, aber darüber
sprachen wir schon.»

Selbst jetzt um die späte Mittagszeit war die Sommerwärme in dieser Höhe
gut zu ertragen. Joh rief zum Essen. Der gedeckte Holztisch vor der
Hütte unter dem vorstehenden Dach wirkte einladend. Alle drei setzten
sich nebeneinander auf die Bank.

Dann sassen sie stumm da. Genossen Johs schmackhafte Zubereitungen und
schauten hinunter ins Tal. Alles war vorläufig gesagt.

Louis wusste, welche Orte einen guten Handyempfang garantierten. Bei
unvorhergesehenen Fragen konnte er Joh oder Manuela anrufen. Auch Mazis
Kontakt hatte er gespeichert. Beim Austauschen ihrer Handynummern war
Louis zusammengezuckt. Er musste erreichbar sein. Und gab gerade ein
Stück Anonymität auf. Konnte ihm die Vernetzung zur Gefahr werden? Er
brauchte einige Sekunden Bedenkzeit, um sich an die Zahlen zu erinnern.
Joh hatte ihn kritisch angesehen, glaubte Louis.

Joh drückte Louis kräftig an sich. Ohne ihn loszulassen. Schob ihn auf
Armeslänge von sich und schaute ihn mit seinem Joh-Blick an.

«Auch wenn wir uns nicht bei dir melden, Ciro, wir sind immer für dich
da. Deine erste Ansprechperson hier oben ist Mazi. Du kannst ihm
bedingungslos vertrauen. Und wie du weisst, führt er dich in alles ein,
was du wissen musst, okay?»

Noch hielt er Louis an den Schultern. Es sah aus, als warte er seine
Reaktion ab. Das Gefühl dieser Nähe stellte sich wieder ein. Und Louis
erinnerte sich, wie Joh ihn gehalten hatte. Wie er in seinen Armen auf
dem Sofa gesessen, sich ergeben, und auch, wie gut es sich angefühlt
hatte.

«Ja, danke, Joh.»

Es war fast so, als führe Joh ihn wie ein Vater in einen neuen
Lebensabschnitt ein. Louis’ Augen waren hinter seiner Sonnenbrille gut
versteckt.

«Ich freue mich auf meine Aufgaben und werde mein Bestes geben. Du
kannst dich auf mich verlassen,» waren Louis’ letzte Worte.

Joh nickte, senkte lässig seine Arme, und Louis wusste, er wollte Joh
nicht enttäuschen.

Solange sie sichtbar blieben, winkten Joh und Manuela ein paar Mal zu
ihm hoch.

Dann verschwanden sie unter dem Licht des späten Nachmittags.

Louis fühlte sich gerade ziemlich aufgeräumt.
